% BHU PYQ -- Manually transcribed & error-corrected
% Paper: GLB-401: Palaeontology and Stratigraphy
% Exam: B.Sc. (Hons.) Semester IV Examination, 2023-24

\documentclass[11pt,a4paper]{article}
\usepackage[a4paper,top=2.5cm,bottom=2.5cm,left=2.8cm,right=2.8cm,headheight=14pt]{geometry}
\usepackage{amsmath,amssymb}
\usepackage[T1]{fontenc}
\usepackage[utf8]{inputenc}
\usepackage{lmodern,microtype,fancyhdr,enumitem,hyperref,lastpage,parskip}

\hypersetup{colorlinks=true,urlcolor=black,linkcolor=black,
  pdftitle={GLB-401 Palaeontology and Stratigraphy -- BHU B.Sc. Sem IV 2023-24}}
\pagestyle{fancy}\fancyhf{}
\renewcommand{\headrulewidth}{0.4pt}\renewcommand{\footrulewidth}{0.4pt}
\fancyhead[L]{\small\textit{Banaras Hindu University}}
\fancyhead[R]{\small\textit{GLB-401: Palaeontology \& Stratigraphy}}
\fancyfoot[C]{\small Page \thepage\ of \pageref{LastPage}}

\newlist{parts}{enumerate}{2}
\setlist[parts,1]{label=(\roman*),leftmargin=*,itemsep=5pt,topsep=3pt}
\newcommand{\pts}[1]{\hfill\textit{[#1]}}

\begin{document}
\begin{center}
  {\large\bfseries BANARAS HINDU UNIVERSITY}\\[2pt]
  {\normalsize B.Sc. (Hons.) --- Semester IV Examination, 2023-24}\\[6pt]
  \rule{0.8\linewidth}{0.5pt}\\[6pt]
  {\large\bfseries Paper: GLB-401 --- Palaeontology and Stratigraphy}\\[4pt]
  {\normalsize Time: 3 Hours \hfill Maximum Marks: 70}\\[2pt]
  {\small Ref. No.: 074445}
\end{center}
\rule{\linewidth}{0.4pt}
\smallskip
\noindent\textit{Instructions: Answer \textbf{five} questions in all, selecting at least \textbf{two} each from Sections A and B, and Section-C which is compulsory. All questions carry equal marks.}
\bigskip

\bigskip\noindent{\large\bfseries SECTION --- A}
\rule{\linewidth}{0.4pt}\smallskip

\noindent\textbf{Question 1.} Write an essay on the subdivisions and scope of Paleontology. \pts{14}

\medskip\noindent\textbf{Question 2.} Describe the morphology of trilobites with the help of neatly labelled diagrams. Comment on their geological history and stratigraphic importance. \pts{14}

\medskip\noindent\textbf{Question 3.} Write in detail the tectonic subdivision of India with suitable illustrations. \pts{14}

\medskip\noindent\textbf{Question 4.} Briefly explain the Vindhyan Supergroup of rocks with special emphasis on lithostratigraphic classification, lithology, and economic importance. \pts{14}

\bigskip\noindent{\large\bfseries SECTION --- B}
\rule{\linewidth}{0.4pt}\smallskip

\noindent\textbf{Question 5.} Write a brief note on \textbf{any two} of the following: \pts{$2 \times 7 = 14$}
\begin{parts}
  \item Define index fossil and discuss its significance.
  \item Differences between regular and irregular Echinoids.
  \item Importance of stratigraphy.
\end{parts}

\medskip\noindent\textbf{Question 6.} Write a brief note on \textbf{any two} of the following: \pts{$2 \times 7 = 14$}
\begin{parts}
  \item Geological distribution of Graptoloidea (Graptoloids).
  \item Order of Superposition in stratigraphy.
  \item Physical subdivision of India with a suitable map.
\end{parts}

\medskip\noindent\textbf{Question 7.} Write in detail on the Precambrian geology of the Dharwar craton with special reference to lithostratigraphic classification, lithology, and economic significance. \pts{14}

\bigskip\noindent{\large\bfseries SECTION --- C}
\rule{\linewidth}{0.4pt}\smallskip

\noindent\textbf{Question 8.} Answer the following in one to five sentences:
\begin{parts}
  \item What are dissepiments? \pts{1}
  \item What is a corona? \pts{1}
  \item What is a nema? \pts{1}
  \item What is a medusa? \pts{1}
  \item What are living fossils? \pts{1}
  \item What are coprolites? \pts{1}
  \item What is an unconformity? \pts{2}
  \item Write the names of two Eons in the Geological Time Scale. \pts{2}
  \item Write the names of two rivers originating from the Himalayas. \pts{2}
  \item Write the names of two mineral deposits located in the Aravalli Craton. \pts{1}
\end{parts}

\vfill\rule{\linewidth}{0.4pt}
\begin{center}\small
\href{https://bhu-exam.vercel.app/}{Click here to visit our website}\\[4pt]
\href{https://me-aryan.vercel.app/}{Click here to visit our contributor}\\[4pt]
\href{https://quashan-me.vercel.app/}{Click here to visit our contributor}
\end{center}
\end{document}
